\documentclass[11pt]{article}
\usepackage{amsmath,amssymb,amsthm,booktabs,mathtools}
\usepackage[margin=1in]{geometry}
\usepackage{microtype}
\usepackage{url}
\newtheorem{theorem}{Theorem}[section]
\newtheorem{lemma}[theorem]{Lemma}
\newtheorem{corollary}[theorem]{Corollary}
\newtheorem{remark}[theorem]{Remark}
\newcommand{\dd}{\,\mathrm d}
\newcommand{\one}{\mathbf 1}
\title{Atlas 130 at High Density: A Rooted-Triangle Moment Plane}
\author{}
\date{}
\begin{document}
\maketitle

\begin{abstract}
Graph Atlas 130 consists of two disjoint triangles joined by one bridge.
For every graphon $W$ of edge density $p\in[2/3,1]$, we prove
\[
 t(H_{130},W)\geq p^3(2p-1)^2,
\]
its exact chromatic target.  The bridge density is written as a quadratic
form in the rooted triangle function.  Passing to the complement kernel and
using $2xy\leq x^2+y^2$ reduces the problem to a single rooted moment
$\int(1-d)\tau^2$.  A pointwise supporting plane in the degree, its one-step
average, and the rooted triangle closes the reduction.  The scalar lemma is
proved by an exact rational conditional Bernstein certificate with 336
boxes and two explicit local corner estimates.  The argument is for
arbitrary graphons; no finite-step reduction is used.  The interval
$1/2<p<2/3$ remains open.
\end{abstract}

\section{The graph and target}

Let $(\Omega,\mu)$ be a probability space and let
$W\colon\Omega^2\to[0,1]$ be symmetric and measurable.  All integrands below
are bounded and nonnegative, so Tonelli--Fubini applies without additional
integrability assumptions.

Atlas 130 has Atlas-labelled edge set
\[
 \{01,02,03,12,34,35,45\},
\]
graph6 string \verb|E{CW|, and degree sequence $(3,3,2,2,2,2)$.
It is isomorphic to two vertex-disjoint triangles joined by an edge between
one chosen vertex of each triangle.  Its chromatic polynomial is
\begin{equation}
 P_{130}(z)=z(z-1)^3(z-2)^2.
 \label{eq:chromatic}
\end{equation}
Indeed, colour the first bridge endpoint in $z$ ways and the other two
vertices of its triangle in $(z-1)(z-2)$ ways.  The second bridge endpoint
then has $z-1$ choices, after which its two triangle partners again have
$(z-1)(z-2)$ ordered choices.  Multiplying gives \eqref{eq:chromatic}, and
also shows $\chi(H_{130})=3$.
Consequently
\begin{equation}
 (1-p)^6P_{130}\!\left(\frac1{1-p}\right)
 =p^3(2p-1)^2.
 \label{eq:target}
\end{equation}
The critical density interval begins at $p=1/2$.

\section{The complement moment reduction}

Write $T_W$ for the integral operator of $W$ and define
\begin{align}
 p&=\int_{\Omega^2}W(x,y)\dd\mu(x)\dd\mu(y),
 &d(x)&=(T_W\one)(x),\notag\\
 a(x)&=(T_Wd)(x),
 &\tau(x)&=\int_{\Omega^2}
 W(x,y)W(x,z)W(y,z)\dd\mu(y)\dd\mu(z).
 \label{eq:rooted}
\end{align}
Let
\[
 T=\int_\Omega\tau(x)\dd\mu(x)=t(K_3,W).
\]
Rooting the two triangles at the bridge endpoints gives
\begin{equation}
 t(H_{130},W)=\int_{\Omega^2}W(x,y)\tau(x)\tau(y)
 \dd\mu(x)\dd\mu(y).
 \label{eq:bridge}
\end{equation}

Put $U=1-W$.  Since $\int\tau=T$, equation \eqref{eq:bridge} becomes
\[
 t(H_{130},W)=T^2-\int_{\Omega^2}U(x,y)\tau(x)\tau(y)\dd\mu^2.
\]
The elementary inequality $2\tau(x)\tau(y)\leq\tau(x)^2+\tau(y)^2$
and symmetry of $U$ imply
\begin{equation}
 \int_{\Omega^2}U(x,y)\tau(x)\tau(y)\dd\mu^2
 \leq\int_\Omega(1-d(x))\tau(x)^2\dd\mu(x).
 \label{eq:moment-reduction}
\end{equation}
Thus
\begin{equation}
 t(H_{130},W)\geq T^2-R,
 \qquad
 R:=\int_\Omega(1-d)\tau^2\dd\mu.
 \label{eq:R}
\end{equation}
This reduction is exact for every balanced complete multipartite graphon.

We shall use the following pointwise feasible region.  It is recorded
explicitly because the extra lower bound in \eqref{eq:feasible} is important
for the high-density supporting plane.

\begin{lemma}[Rooted feasible region]\label{lem:feasible}
For almost every $x$, the tuple $(d,a,t)=(d(x),a(x),\tau(x))$ satisfies
\begin{equation}
\begin{gathered}
 0\leq d\leq1,
 \qquad \max\{0,d+p-1\}\leq a\leq d,\\
 \max\{0,2a-p,a-d(1-d)\}\leq t\leq\min\{a,d^2\}.
\end{gathered}
\label{eq:feasible}
\end{equation}
Moreover,
\begin{equation}
 \int d\dd\mu=p,
 \qquad \int a\dd\mu=\int d^2\dd\mu,
 \qquad \int t\dd\mu=T.
 \label{eq:means}
\end{equation}
\end{lemma}

\begin{proof}
The upper bounds $a\leq d$, $t\leq a$, and $t\leq d^2$ follow by deleting
one or more $[0,1]$-valued factors from the defining integrals.  Also
\[
 a=p-\int_\Omega(1-W(x,y))d(y)\dd\mu(y)
 \geq p-(1-d)=d+p-1.
\]
For the rooted Goodman bound, apply
$W(x,y)W(x,z)\geq W(x,y)+W(x,z)-1$, multiply by $W(y,z)$,
and integrate to get $t\geq2a-p$.  Applying instead
$W(x,z)W(y,z)\geq W(x,z)+W(y,z)-1$, multiplying by $W(x,y)$,
and integrating gives
\[
 t\geq\int W(x,y)(d(x)+d(y)-1)\dd\mu(y)
 =a-d(1-d).
\]
Nonnegativity supplies the remaining lower bounds.  The three identities in
\eqref{eq:means} are symmetry and Fubini.
\end{proof}

\section{A full-interval same-orientation companion}

The following sharp inequality holds already on the complete chromatic
interval.  It does not by itself settle Atlas 130, but it isolates the sole
remaining crossed-edge issue without losing any rooted-triangle surplus.

For $d=0$, the nonnegative function $W(x,\cdot)$ vanishes almost everywhere,
so $a=\tau=0$ there.  We may therefore use the convention $a^2/d=0$ on
$\{d=0\}$ and put
\[
 \mathcal R=\int_\Omega\frac{a^2}{d}\dd\mu,
 \qquad F=\int_\Omega a\tau\dd\mu,
 \qquad S=\int_\Omega d\tau^2\dd\mu.
\]

\begin{lemma}[Strengthened $T_Wd$-weighted triangle]
\label{lem:strengthened-weight}
For every graphon of density $p\in[1/2,1]$,
\begin{equation}
 F\geq(2p-1)\mathcal R.
 \label{eq:strengthened-weight}
\end{equation}
\end{lemma}

\begin{proof}
Write $M_j=\int d^j$, $c=2p-1$, $q=1-p$, and
$V=M_2-p^2$.  The rooted Goodman inequality in
Lemma~\ref{lem:feasible}, multiplied by $a\geq0$, gives
\begin{equation}
 F\geq2\int a^2-pM_2\geq2M_2^2-pM_2,
 \label{eq:F-moment}
\end{equation}
because $\int a=M_2$ and Jensen applies to $a^2$.

Let $\mu_3=\int(d-p)^3$.  Exact expansion gives
\begin{equation}
 2M_2^2-pM_2-cM_3
 =2V^2+2pqV-c\mu_3.
 \label{eq:third-moment-identity}
\end{equation}
Pointwise,
\[
 (d-p)^3\leq q(d-p)^2,
\]
since their difference in the other direction is
$(d-p)^2(1-d)\geq0$.  Hence $\mu_3\leq qV$.  As $c\geq0$, the right side of
\eqref{eq:third-moment-identity} is at least
\[
 2V^2+2pqV-cqV=V(2V+q)\geq0.
\]
Thus \eqref{eq:F-moment} yields $F\geq cM_3$.

For $d(x)>0$, the ratio $a(x)/d(x)$ is the average of $d(y)$ under the
neighbour probability measure $W(x,y)\dd\mu(y)/d(x)$.  Conditional Jensen
therefore gives
\[
 \frac{a(x)^2}{d(x)}
 \leq\int_\Omega W(x,y)d(y)^2\dd\mu(y).
\]
After integration and symmetry, $\mathcal R\leq M_3$.  Multiplication by
$c\geq0$ completes \eqref{eq:strengthened-weight}.
\end{proof}

\begin{corollary}[Sharp same-orientation amalgam]
\label{cor:same-orientation}
For $p\in[1/2,1]$,
\begin{equation}
 S\geq cF\geq c^2\mathcal R\geq p^3c^2.
 \label{eq:same-orientation-chain}
\end{equation}
The left side is the density of Graph Atlas 117 (graph6
\verb|EtaG|), two triangles sharing their root with one leaf at that root.
\end{corollary}

\begin{proof}
Weighted Cauchy--Schwarz gives
\[
 F^2
 =\left(\int (\sqrt d\,\tau)\frac a{\sqrt d}\right)^2
 \leq S\mathcal R.
\]
If $\mathcal R=0$, then $F=0$ and the first inequality is immediate.
Otherwise Lemma~\ref{lem:strengthened-weight} gives
$S\geq F^2/\mathcal R\geq cF$.
Another weighted Cauchy--Schwarz inequality and $M_2\geq p^2$ give
\[
 M_2^2=\left(\int a\right)^2\leq p\mathcal R,
 \qquad \mathcal R\geq p^3.
\]
Combining these estimates proves \eqref{eq:same-orientation-chain}.  Rooting
the two triangles and the leaf at their common vertex gives the asserted
Atlas 117 density identity.  Its chromatic polynomial is
$z(z-1)^3(z-2)^2$, so the last member of the chain is its exact target.
\end{proof}

For Atlas 130 the analogous expression is crossed across the bridge:
\[
 t(H_{130},W)=\int W(x,y)\tau(x)\tau(y)\dd\mu^2,
 \qquad
 F=\int W(x,y)\tau(x)d(y)\dd\mu^2.
\]
Consequently the single comparison
\begin{equation}
 t(H_{130},W)\stackrel{?}{\geq}cF
 \label{eq:remaining-comparison}
\end{equation}
would combine with \eqref{eq:same-orientation-chain} to prove Atlas 130 on
$[1/2,1]$.  Corollary~\ref{cor:same-orientation} proves the corresponding
same-orientation comparison $S\geq cF$; it does not identify the two edge
orientations, and no such identification is assumed below.

\begin{remark}[Exact boundary of the scalar moment relaxation]
\label{rem:low-scalar-obstruction}
The missing interval cannot be recovered merely by combining
\eqref{eq:moment-reduction} with the pointwise region
\eqref{eq:feasible} and the mean identities \eqref{eq:means}.  At
$p=3/5$, consider a two-atom probability distribution with weights
$17/35,18/35$ on the rooted tuples
\[
 (d,a,t)=
 \left(\frac{21}{34},\frac{5673}{11560},\frac{441}{1156}\right),
 \qquad
 \left(\frac7{12},\frac{1451}{6120},0\right).
\]
Every one of the eight polynomial slacks used in
\eqref{eq:feasible} is nonnegative at both atoms.  Direct calculation gives
\[
 \mathbb E d=\frac35,
 \qquad
 \mathbb E a=\mathbb E d^2=\frac{49}{136},
 \qquad
 T=\mathbb E t=\frac{63}{340},
\]
and
\[
 R=\mathbb E[(1-d)t^2]=\frac{361179}{13363360}.
\]
Nevertheless,
\[
 T^2-R-p^3(2p-1)^2
 =-\frac{11138769}{8352100000}<0.
\]
This distribution is not asserted to come from a graphon, so it is not a
counterexample to Atlas 130.  It instead shows exactly why a low-density
proof must use additional kernel compatibility, such as the crossed
comparison \eqref{eq:remaining-comparison}; no different supporting plane
on the presently recorded scalar feasible region can suffice.
\end{remark}

\section{The high-density supporting plane}

Set
\begin{equation}
 c=2p-1,qquad m=pc,qquad A=p^2c^2(1-p),
 \label{eq:coefficients1}
\end{equation}
and define
\begin{equation}
 \gamma=\frac{-9p^3-p^2+5p-1}{2},
 \qquad
 \beta=\gamma-5p^2c^2.
 \label{eq:coefficients2}
\end{equation}

\begin{lemma}[Rooted moment plane]\label{lem:plane}
Let $2/3\leq p\leq1$.  Every $(d,a,t)$ satisfying
\eqref{eq:feasible} obeys
\begin{equation}
 (1-d)t^2
 \leq A+2m(t-m)+\beta(d-p)+\gamma(a-d^2).
 \label{eq:plane}
\end{equation}
\end{lemma}

\begin{proof}
Put the right side minus the left side of \eqref{eq:plane} equal to
$E_p(d,a,t)$.  Direct Bernstein coefficients show that $\gamma<0$ throughout
$[2/3,1]$.  For fixed $p,d,t$, the minimum of $E_p$ over feasible $a$ is
therefore attained at
\[
 a_{\max}=\min\left\{d,\frac{t+p}{2},t+d(1-d)\right\}.
\]
Call these three faces A, B, and C.  On face A feasibility forces
$(a,t)=(d,d^2)$.  On face B the residual is concave in $t$, and its possible
interval endpoints are
\[
\begin{array}{c|c|c}
 &a&t\\ \hline
 \mathrm{B0}&p/2&0\\
 \mathrm{BL1}&d^2-d+p&p-2d+2d^2\\
 \mathrm{BL2}&d+p-1&2d+p-2\\
 \mathrm{BU1}&(d^2+p)/2&d^2\\
 \mathrm{BU2}&d&2d-p\\
 \mathrm{BUp}&p&p.
\end{array}
\]
On face C the lower endpoints are
\[
 (a,t)=(d(1-d),0),\qquad(d+p-1,d^2+p-1),
\]
called C0 and CL; its upper endpoints duplicate A and BL1.  Concavity in
$t$ therefore reduces the lemma to the nine rows just listed.

For completeness, here is the exact certificate.  Substitute
\[
 p=\frac{2+X}{3},\qquad d=D,qquad (X,D)\in[0,1]^2.
\]
At every endpoint impose the eight polynomial feasibility constraints
\begin{equation}
 a,\ d-a,\ a-d-p+1,\ t,\ t-2a+p,\
 t-a+d(1-d),\ a-t,\ d^2-t\ \geq0.
 \label{eq:eight}
\end{equation}
Each target and constraint is converted exactly to a common bidegree
Bernstein tensor.  A box is discarded if one constraint has all coefficients
strictly negative, and accepted if all target coefficients are nonnegative;
otherwise its widest side is bisected by exact integer de Casteljau
subdivision.  The resulting tree counts are
\begin{center}
\begin{tabular}{lrrr}
\toprule
face & visited & accepted & infeasible\\
\midrule
A&81&19&22\\
B0&7&0&4\\
BL1&9&4&1\\
BU1&25&12&1\\
BU2&35&9&9\\
BUp&17&4&5\\
C0&17&5&4\\
CL&145&73&0\\
\midrule
total&336&126&46\\
\bottomrule
\end{tabular}
\end{center}
The BL1 residual is divided first by its manifest nonnegative contact square:
\begin{align}
 E_{\mathrm{BL1}}
 &=4(d-p)^2P(p,d),\notag\\
 P(p,d)&=d^3+2d^2p-3d^2+3dp^2-5dp+3d
 +4p^3-5p^2+3p-1.
 \label{eq:BL1}
\end{align}
The conditional certificate is applied to $P$.  On BL2, one of
\eqref{eq:eight} is $-(D-1)^2\geq0$, so feasibility forces $D=1$; the
remaining univariate residual is
\begin{equation}
 \frac8{243}(X-1)^2(X+2)^2(2X+1)\geq0.
 \label{eq:BL2}
\end{equation}

Only the corner $(X,D)=(1,1)$ needs local treatment on A and CL.  Put
$r=1-X$ and $u=1-D$.  On A, feasibility contains
$u^2-r/3\geq0$.  For $0\leq r,u\leq\varepsilon_A=1/64$, exact expansion gives
\[
 E_A=9u^2-\frac r3+R_A.
\]
If $R_A=\sum c_{ij}r^iu^j$, every term either has $i\geq1$ or has
$i=0,j\geq3$, and exact arithmetic gives
\begin{equation}
 \sum_{i\geq1}3|c_{ij}|\varepsilon_A^{i-1+j}
 +\sum_{j\geq3}|c_{0j}|\varepsilon_A^{j-2}
 =\frac{1247857961}{2717908992}<8.
 \label{eq:A-loss}
\end{equation}
Since $r\leq3u^2$, this bounds $|R_A|$ by the displayed constant times
$u^2$, whereas $9u^2-r/3\geq8u^2$.

On CL the quadratic part is
\[
 Q(r,u)=\frac89r^2-\frac{16}{3}ru+9u^2.
\]
After subtracting $\frac1{20}(r^2+u^2)$, its leading matrix has first
principal minor $151/180$ and determinant $1429/3600$, so
$Q\geq\frac1{20}(r^2+u^2)$.  With
$\varepsilon_C=1/512$, all terms of $R_C:=E_{CL}-Q$ have total degree at
least three and
\begin{equation}
 \sum |c_{ij}|\varepsilon_C^{i+j-2}
 =\frac{2644746959}{65229815808}<\frac1{20}.
 \label{eq:C-loss}
\end{equation}
Thus $|R_C|<\frac1{20}(r^2+u^2)$ on that corner.  These two local estimates
close precisely the boxes exempted from subdivision.  Every operation in
the certificate is rational, completing the proof.
\end{proof}

\section{The arbitrary-graphon theorem}

\begin{theorem}\label{thm:main}
For every graphon $W$ of edge density $p\in[2/3,1]$,
\[
 \boxed{t(H_{130},W)\geq p^3(2p-1)^2.}
\]
\end{theorem}

\begin{proof}
Apply Lemma~\ref{lem:plane} pointwise to
$(d(x),a(x),\tau(x))$ and integrate.  The two centered terms disappear by
\eqref{eq:means}, giving
\begin{equation}
 R\leq p^2c^2(1-p)+2pc(T-pc).
 \label{eq:R-bound}
\end{equation}
Combining \eqref{eq:R} and \eqref{eq:R-bound} yields
\begin{align*}
 t(H_{130},W)-p^3c^2
 &\geq T^2-p^2c^2(1-p)-2pc(T-pc)-p^3c^2\\
 &=(T-pc)^2\geq0.
\end{align*}
This last step is an identity.  In particular, although Goodman's inequality
$T\geq pc$ is compatible with the proof, its direction is not needed after
the square completion.
\end{proof}

\section{Scope and reproducibility}

The theorem covers the whole high subinterval $[2/3,1]$ for arbitrary
graphons.  It does not settle $1/2<p<2/3$.  The same complement moment
reduction \eqref{eq:moment-reduction} remains valid there, but the scalar
support problem needs additional global information; the high-density plane
is not asserted outside its certified interval.

No earlier graph-density theorem is used: the analytic inputs are only
Tonelli--Fubini, the union bound on two numbers in $[0,1]$, elementary
Cauchy--Schwarz and Jensen inequalities, the third-moment calculation
\eqref{eq:third-moment-identity}, and $2xy\leq x^2+y^2$.  At $p=1$, the usual identity
$\int(1-W)=0$ forces $W=1$ almost everywhere and both sides of the theorem
equal one; the displayed proof and exact scalar certificate also include
this endpoint without division.  At $p=2/3$ the same formulas apply
directly.  Vanishing degrees cause no difficulty because no quotient by
$d$, $T$, $p-1$, or $2p-1$ occurs.

The script
\begin{center}
\texttt{codes/verify\_atlas130\_high\_density\_root\_moment.py}
\end{center}
checks the Atlas identification, graph6 string, chromatic polynomial and
target, all symbolic face reductions, the sign of $\gamma$, the 336-box
conditional Bernstein certificate, the BL1 and BL2 factorizations, and both
local remainder estimates using exact rational arithmetic.  It also performs
25 exact two-block sanity checks of the rooted bridge and complement-moment
identities.  No claim of Lean verification is made.

\end{document}
